\chapter{从约束到Lagrange方程}

\section{单摆为什么只需要一个方程}
平面单摆的质点有两个笛卡尔坐标，却必须满足$x^2+y^2=l^2$。
直接使用Newton方程，要同时求$x,y$和未知张力；用摆角$\theta$描述位置，约束自动满足，
张力又与允许的瞬时位移垂直，因此可以在求摆角运动时消去。
Lagrange方法把这一想法推广到许多自由度。

设$N$个粒子受到独立的完整约束$f_\alpha(\bm r_1,\ldots,\bm r_N,t)=0$，
局部可以用$n$个独立广义坐标表示位置：$\bm r_a=\bm r_a(q_1,\ldots,q_n,t)$。
广义坐标可以是长度、角度或其他参数，未必具有相同量纲。
“独立”要求约束的Jacobian具有所需的满秩；坐标在极点等位置退化时应换坐标。

在同一时刻比较两条允许位置，得到虚位移
\begin{equation}
 \delta\bm r_a=\sum_i\pv{\bm r_a}{q_i}\delta q_i,\qquad \delta t=0.
\end{equation}
它不同于实际位移$\d\bm r_a=\sum_i(\p\bm r_a/\p q_i)\d q_i+(\p\bm r_a/\p t)\d t$。
例如运动支座的约束力可对实际运动做功，同时对允许虚位移不做功。

\section{从Newton方程推导}
以下在惯性系中使用Newton方程，并取各粒子质量$m_a$恒定。
把力分为主动力$\bm F_a$与约束力$\bm R_a$。
对于理想约束，$\sum_a\bm R_a\cdot\delta\bm r_a=0$，于是
\begin{equation}
 \sum_a(\bm F_a-m_a\ddot{\bm r}_a)\cdot\delta\bm r_a=0.
 \label{eq:cm-dalembert}
\end{equation}
这就是d'Alembert形式。定义$T=\tfrac12\sum_a m_a\dot{\bm r}_a^2$与广义力
$Q_i=\sum_a\bm F_a\cdot\p\bm r_a/\p q_i$。
由$\p\dot{\bm r}_a/\p\dot q_i=\p\bm r_a/\p q_i$及链式法则，
\begin{align}
 \dv{}{t}\pv{T}{\dot q_i}
 &=\sum_a m_a\ddot{\bm r}_a\cdot\pv{\bm r_a}{q_i}
 +\sum_a m_a\dot{\bm r}_a\cdot\dv{}{t}\pv{\bm r_a}{q_i},\\
 \pv T{q_i}&=\sum_a m_a\dot{\bm r}_a\cdot\pv{\dot{\bm r}_a}{q_i}
 =\sum_a m_a\dot{\bm r}_a\cdot\dv{}{t}\pv{\bm r_a}{q_i}.
\end{align}
相减并利用独立的$\delta q_i$，得到
\begin{equation}
 \dv{}{t}\pv T{\dot q_i}-\pv T{q_i}=Q_i.
\end{equation}
若主动力可分为$-\nabla_aV$与其余部分，令$L=T-V$，则
\begin{equation}
 \boxed{\dv{}{t}\pv L{\dot q_i}-\pv L{q_i}=Q_i^{\rm nc}.}
 \label{eq:cm-el}
\end{equation}
广义力与对应坐标互为功的共轭：$Q_i\delta q_i$的量纲为能量。
角坐标对应的广义力通常是力矩。

式\eqref{eq:cm-el}适用于固定或显含时间的理想完整约束。
对不可积的速度约束（非完整约束），须从允许虚位移与约束力出发，通常不能直接消元再对作用量变分；本章只对完整约束这样消元。

\begin{example}{移动支点的单摆}{cm-moving-pendulum}
支点沿水平方向按给定函数$a(t)$运动，取
$x=a(t)+l\sin\theta$，$y=-l\cos\theta$。
忽略仅依赖时间的加项后，
\begin{equation}
 L=\frac12ml^2\dot\theta^2+ml\dot a\dot\theta\cos\theta+mgl\cos\theta.
\end{equation}
Euler--Lagrange方程给出
\begin{equation}
 \ddot\theta+\frac gl\sin\theta=-\frac{\ddot a}{l}\cos\theta.
\end{equation}
支点的加速度而非速度驱动摆角；$\ddot a=0$时恢复普通单摆，符合Galileo相对性。
该例中支点可以做功，不能预先假定摆的机械能守恒。
\end{example}

\section{驻定作用量，而不总是最小作用量}
对于无非保守广义力的系统，定义
\begin{equation}
 S[q]=\int_{t_1}^{t_2}L(q,\dot q,t)\d t.
\end{equation}
取$q_i(t)\mapsto q_i(t)+\epsilon\eta_i(t)$，其中$\eta_i(t_1)=\eta_i(t_2)=0$。
分部积分给出
\begin{align}
 \delta S
 &=\sum_i\Bigb{\pv L{\dot q_i}\delta q_i}_{t_1}^{t_2}
 +\int_{t_1}^{t_2}\sum_i\Bigp{\pv L{q_i}-\dv{}{t}\pv L{\dot q_i}}\delta q_i\d t.
\end{align}
端点项为零，对任意内部变分令$\delta S=0$就得到Lagrange方程。
这固定了两端的时刻与位置，并不要求粒子预先“知道终点”；预测仍可用初值求解。
驻定路径可能是局部最小或鞍点，不能一概称为“最小”。

若$L'=L+\d F(q,t)/\d t$，作用量只增加端点差，故固定端点方程不变。
这解释了为什么不同形式的Lagrangian可以描述同一运动。

\section{约束反力与耗散怎样放回来}
若仍希望求约束力，可保留冗余坐标，并对完整约束$f_\alpha(q,t)=0$引入乘子：
\begin{equation}
 \dv{}{t}\pv L{\dot q_i}-\pv L{q_i}
 =Q_i^{\rm nc}+\sum_\alpha\lambda_\alpha\pv{f_\alpha}{q_i}.
\end{equation}
配合约束方程一起求解。$\lambda_\alpha$的数值和量纲依赖$f_\alpha$的归一化，
真实约束力是其与约束梯度的乘积；不要未经换算就把乘子本身叫作张力。

对线性黏滞阻尼，若广义力可写为$Q_i^{\rm nc}=-\sum_jc_{ij}\dot q_j$，
其中$C$对称半正定，可定义Rayleigh耗散函数
\begin{equation}
 \mathcal R=\frac12\sum_{ij}c_{ij}\dot q_i\dot q_j,\qquad
 Q_i^{\rm nc}=-\pv{\mathcal R}{\dot q_i},\qquad
 \sum_iQ_i^{\rm nc}\dot q_i=-2\mathcal R\leq0.
\end{equation}
$\mathcal R$具有功率量纲，不是附加势能；不能简单写$L=T-V-\mathcal R$后使用无耗散方程。
干摩擦等力不一定具有这样的二次耗散函数。

\section{循环坐标、能量与Noether关系}
定义广义动量$p_i=\p L/\p\dot q_i$。当$\p L/\p q_j=0$且$Q_j^{\rm nc}=0$时，
$q_j$称为循环坐标，$p_j$守恒。
例如中心势中$\phi$循环，$p_\phi=\mu r^2\dot\phi$就是角动量。

定义能量函数$E_L=\sum_i\dot q_ip_i-L$，沿运动方程有
\begin{equation}
 \dv{E_L}{t}=-\pv L t+\sum_iQ_i^{\rm nc}\dot q_i.
 \label{eq:cm-lagrange-energy}
\end{equation}
所以无非保守力且$\p L/\p t=0$时$E_L$守恒。
当$T$对$\dot q$为二次齐次式且$V$与速度无关时，由Euler齐次函数定理直接得到$E_L=T+V$。
显含时间的坐标变换或速度相关相互作用，需要实际计算后再判断能量的物理含义。

对不改变时间的无穷小变换$\delta q_i=\epsilon\xi_i(q,t)$，若对任意允许路径
（而非仅对某条运动方程的解）都有$\delta L=\epsilon\d B/\d t$，由运动方程可得
\begin{equation}
 \delta L=\epsilon\dv{}{t}\biggp{\sum_i p_i\xi_i},\qquad
 \mathcal Q=\sum_i p_i\xi_i-B=\const.
\end{equation}
这是一种常用的Noether关系：连续对称性可以允许Lagrangian改变一个全导数。
空间平移给出动量，旋转给出角动量；时间平移的结论由式\eqref{eq:cm-lagrange-energy}体现。
上述推导对在壳运动且无非保守广义力成立，不能把有外部驱动的粒子子系统任意视为封闭系统。

\begin{example}{匀速旋转竖直圆环上的小珠}{cm-hoop}
半径$a$的光滑竖直圆环绕竖直直径以给定角速度$\Omega$转动，
用从最低点量起的$\theta$描述小珠。惯性系动能和势能给出
\begin{equation}
 L=\frac12ma^2\bigl(\dot\theta^2+\Omega^2\sin^2\theta\bigr)+mga\cos\theta,
\end{equation}
从而
\begin{equation}
 \ddot\theta=\sin\theta\Bigp{\Omega^2\cos\theta-\frac ga},\qquad
 V_{\rm eff}=-mga\cos\theta-\frac12ma^2\Omega^2\sin^2\theta.
\end{equation}
最低点在$\Omega^2<g/a$时稳定，小振动频率为$\sqrt{g/a-\Omega^2}$。
当$\Omega^2>g/a$时，出现$\cos\theta_0=g/(a\Omega^2)$的两个稳定平衡，
其小振动频率为$|\Omega\sin\theta_0|$；最低点变得不稳定。
临界点$\Omega^2=g/a$的二阶项消失，但四阶项为正，仍需非线性分析。
这里守恒的是$E_L=ma^2\dot\theta^2/2+V_{\rm eff}$，
不是小珠的惯性系机械能；维持圆环匀速转动的电机参与能量交换。
\end{example}

